% ===========================================================================
% Abstract (under 200 words, the course's rule; 199 now, so a sentence added needs one removed):
% why, method, results, what is new. tests/test_numbers_full.py counts the words.

% ===========================================================================
\begin{abstract}
\noindent
Literature assistants such as SciSpace and Consensus answer scientific questions with
citations, but each answer depends on one step they do not expose: which is finding the right
evidence. This paper takes only that step and improves it. A small sentence encoder is
fine-tuned on scientific claims and is fused with keyword search (BM25) at a weight chosen on
held-out claims, and the cross-encoder reranker of the standard hybrid pipeline is also removed.
We prove, that below a BM25 weight of $61/160 \approx 0.38$, reciprocal rank fusion can only
re-order the dense retriever's candidates, this also explains why the chosen weight works. On
SciFact (300 test claims with known evidence) the improved retrieval reaches an nDCG@10 of
0.759 against 0.715 for the standard pipeline ($p = 0.003$), which is at about a third of the time per
claim. On a second dataset, NFCorpus, it is statistically level with the reranked pipeline
(0.350 against 0.358, $p = 0.35$) at a sixth of the time, while fine-tuning helps less on its
lay queries. A failure analysis performed shows that 42\% of the hardest misses are evidence the
benchmark marks but that does not state the claim, and that the improvement repairs 22 of the
40 hardest failures.

\medskip
\noindent\textbf{Keywords:} evidence retrieval; scientific claim verification; hybrid
retrieval; reciprocal rank fusion; dense retrieval; fine-tuning; cross-encoder reranking
\end{abstract}
